A telescope with an achromatic convex objective lens of diameter
$D = 15$\,cm and focal length $f = 200$\,cm is pointed to a star at the
zenith.

\begin{description}
\item[(T09.1)] Find the diameter (in m), $d_{\mathrm{image}}$, of the image
  of a point source as produced by the objective lens at its focal plane for
  green light ($\lambda = 550$\,nm), considering only the effects of
  diffraction. \hfill[1]
\end{description}

The image of an astronomical source is also affected by the so-called
``atmospheric seeing''. The boundaries between the layers in the atmosphere
as well as the refractive indices of the layers change continuously due to
turbulence, temperature variation and other factors. This leads to tiny
changes in the position of the image in the focal plane of the telescope,
known as the ``twinkling effect''. For rest of the problem, apart from the
diffraction limited finite size of the image of the star discussed above, no
interference effects will be considered.

The left panel of the figure below shows a vertical cross-section of the
atmosphere with multiple layers of different refractive indices
($n_1, n_2, n_3, \ldots$). The right panel shows the zoomed in view of a
thin vertical segment of the atmosphere and the boundary between the two
lowest atmospheric layers of refractive indices $n_1$ and $n_2$
($n_1 > n_2$). We consider only these two layers and their boundary for this
problem. The diagrams are not to scale.

\ioaafig{2025_TH_T09-f1.pdf}

\begin{description}
\item[(T09.2)] Let the boundary between the two layers be at a height
  $H = 1$\,km directly above the telescope objective, with a tilt of
  $\theta = 30^\circ$ with respect to the horizontal plane. In all parts of
  this problem $\theta$ is taken to be positive in the anti-clockwise
  direction. For a monochromatic light source, $n_1 = 1.00027$ and
  $n_2 = 1.00026$. Let the angular shift in the position of the image at the
  focal plane of the telescope for a star at the zenith be $\alpha$.
  \begin{description}
  \item[(T09.2a)] Draw an appropriately labelled ray-diagram at the boundary
    showing $n_1$, $n_2$, $\theta$ and $\alpha$. \hfill[2]
  \item[(T09.2b)] Find the expression for $\alpha$ in terms of $\theta$,
    $n_1$ and $n_2$. Use the small angle approximations:
    $\sin\alpha \approx \alpha$ and $\cos\alpha \approx 1$. \hfill[2]
  \item[(T09.2c)] Calculate the displacement, $\Delta x_\theta$ (in m), in
    the position of the image if $\theta$ increases by 1\% (keeping $n_1$
    and $n_2$ fixed). \hfill[3]
  \item[(T09.2d)] Calculate the displacement, $\Delta x_n$ (in m), in the
    position of the image if $n_2$ increases by 0.0001\% (keeping $n_1$ and
    $\theta$ fixed). \hfill[3]
  \end{description}
\item[(T09.3)] For white light coming from a star at the zenith, choose
  which of the following most closely describes the shape and colour of the
  image by ticking the appropriate box (only one) in the Summary
  Answersheet. Note $x$ increases from left to right in the figure.
  \hfill[2]

  \begin{tabular}{clccc}
  & Image colour & Image shape & Left edge & Right edge \\
  A & White & Circular & & \\
  B & White & Elliptical & & \\
  C & Coloured & Circular & Blue & Red \\
  D & Coloured & Circular & Red & Blue \\
  E & Coloured & Elliptical & Blue & Red \\
  F & Coloured & Elliptical & Red & Blue \\
  \end{tabular}
\end{description}

For all remaining parts of this question we consider monochromatic green
light with $\lambda = 550$\,nm. We model the boundary between the layers as
a set of infinite zigzag planes (running perpendicular to the plane of the
page) separated by $d = 10$\,cm along $x$-axis, with either
$\theta = 10^\circ$ or $\theta = -10^\circ$. The figure below (not to scale)
shows a cross-section of this model of the atmosphere of width $W$
($W \ll H$). For telescopes with large aperture this zigzag nature of the
boundary results in formation of speckles in the focal plane.

\ioaafig{2025_TH_T09-f2.pdf}

\begin{description}
\item[(T09.4)] Consider an atmosphere modelled as above.
  \begin{description}
  \item[(T09.4a)] A section of the atmosphere with consecutive zigzag
    planes, with same parameters as stated above, is shown in the diagram
    below (not to scale).

    \ioaafig{2025_TH_T09-f3.pdf}

    In this diagram, reproduced in the Summary Answersheet, draw the paths
    of the incident light rays up to the plane where the telescope objective
    is placed, shown by the gray dotted line. Mark the region(s), if any, by
    ``X'' in the diagram where no light rays will reach. \hfill[4]
  \item[(T09.4b)] Calculate the width $W_X$ of such region(s). \hfill[3]
  \item[(T09.4c)] Find the largest diameter, $D_{\max}$, of the telescope
    objective with which it will be possible to obtain a single image of a
    star, by appropriately choosing the location of the telescope relative
    to the structure of the boundary. \hfill[4]
  \end{description}
\item[(T09.5)] Consider the case when the zigzag shape of the boundary is
  allowed in both $x$ and $y$ directions, (like a field of pyramids), and
  $D = 100$\,cm (with $f = 200$\,cm). Draw the qualitative pattern of the
  resulting speckles in the box given in the Summary Answersheet. \hfill[6]
\item[(T09.6)] For a turbulent atmosphere again consider the same parallelly
  running zigzag shape of the boundary layer only along $x$-direction, but
  now the angle between two planes are changing at a uniform rate from
  $10^\circ$ to $-10^\circ$ in 1.0\,s. Assume that this leads to a uniform
  rate of shift of the position of the image.

  Consider a telescope with $D = 8$\,cm and $f = 1$\,m. Estimate the longest
  exposure time $t_{\max}$ allowed for its CCD camera so that one gets only
  a single image, and any possible deviation in its position remains less
  than 1\% of the diffraction limited diameter of the image. \hfill[5]
\end{description}
